%%%PAGE 11 rot=0 legibility=good summary=静电场：三点电荷平衡比例、场强/电势/电场力做功公式、微观粒子是否计重力、电势几何法（匀强与非匀强）、等量电荷中垂线电场与E_max
%%%BLOCK id=011-01 ch=09 sec=库仑定律·三电荷平衡 kind=formula alt=none cont=prev y=0.04-0.14
\[
q_1:q_2:q_3=L_1^2(L_1+L_2)^2:L_1^2L_2^2:L_2^2(L_1+L_2)^2
\]
% 以下三组口诀依次写在上式三项的正下方
邻邻总总\quad 左左右右\quad \unclear{林林总总}
%FIG p011_a 0.20 0.093 0.49 0.14
\notefig[0.3]{figures/p011_a.png}
%%%END
%%%BLOCK id=011-02 ch=09 sec=场强·电势·电场力做功 kind=formula alt=none cont=none y=0.14-0.25
\[
E=\underbrace{\frac{F}{q}}_{\text{定义式}}=\underbrace{\frac{kQ}{r^2}}_{\substack{\text{点电荷}\\\text{决定式}}}=\underbrace{\frac{U}{d}}_{\text{匀强}}
\]
\[
\varphi=\frac{E_p}{q}=\frac{kq}{r}\quad\mbox{\unclear{带正负}}
\] % CHECK: 手写分式似为 E/d（分母像 d），按物理含义转为 E_p/q
\[
qEs\cos\theta=W_{ab}=qU_{ab}=q(\varphi_a-\varphi_b)=E_{pa}-E_{pb}
\]
%%%END
%%%BLOCK id=011-03 ch=09 sec=带电粒子·是否计重力 kind=note alt=none cont=none y=0.25-0.31
微观粒子（质中电）无重力\qquad 油滴小球微粒 有重力
%%%END
%%%BLOCK id=011-04 ch=09 sec=电势·几何法 kind=method alt=none cont=none y=0.33-0.60
\textbf{电势几何}\quad 匀强电场\unclear{电场线}必平行且等距

\textbf{匀强电场}：
\begin{itemize}
\item 平行四边形，对角线电势和相等
%FIG p011_b 0.478 0.388 0.566 0.487
\notefig[0.25]{figures/p011_b.png}
\[
\varphi_A-\varphi_B=\varphi_D-\varphi_C,\qquad \varphi_A+\varphi_C=\varphi_B+\varphi_D
\]
\item 三个点电势 $\to E$

先找等势面 $\perp E$ $\to E$方向

$E=\dfrac{U}{d}$
\item % 手写大括号只括到上面几项，此行在括号之外
若 B 为 AC 中点，$\varphi_B=\dfrac{\varphi_A+\varphi_C}{2}$
\end{itemize}
%%%END
%%%BLOCK id=011-05 ch=09 sec=电势·几何法 kind=method alt=none cont=none y=0.60-0.745
\textbf{非匀强电场}
%FIG p011_c 0.245 0.605 0.375 0.73
\notefig[0.25]{figures/p011_c.png}
\begin{align*}
&\varphi_A>\varphi_B>\varphi_C\qquad\text{匀强}\ \varphi_B=\dfrac{\varphi_A+\varphi_C}{2}\\
&U_{AB}=\overset{\text{大}}{\overline{E}_{AB}}\cdot|AB|\\
&U_{BC}=\overset{\text{小}}{\overline{E}_{BC}}\cdot|BC|\qquad U_{AB}>U_{BC}\\
&\varphi_A-\varphi_B>\varphi_B-\varphi_C\\
&\therefore\ \varphi_B<\tfrac{1}{2}(\varphi_A+\varphi_C)
\end{align*}
%%%END
%%%BLOCK id=011-06 ch=09 sec=等量电荷·中垂线电场 kind=knowledge alt=none cont=none y=0.745-1.0
%FIG p011_d 0.10 0.74 0.284 0.858
\notefig[0.25]{figures/p011_d.png}
中垂\unclear{线等}势面：等量异电

中垂线\unclear{/}两条电场线：等量同电

\begin{keybox}
\red{中轴线上 $\theta=35.6^\circ$ 时 $E_{\max}$}\quad\red{后有证明} % CHECK: tan θ=√2/2 对应 θ≈35.3°，手写为 35.6°
\end{keybox}

\[
\varphi_A>\varphi_B>\varphi_C\qquad \tan\theta=\frac{\sqrt{2}}{2}\ \text{处，}
\]
\[
E_{c\max}>E_A>E_B
\]
%FIG p011_e 0.085 0.893 0.35 1.0
\notefig[0.26]{figures/p011_e.png}
中垂线等势面（$\varphi=0$）
%%%END
%%%PAGEEND 11
